\documentclass[11pt]{article}
\usepackage[a4paper,margin=22mm]{geometry}
\usepackage{fontspec}
\setmainfont{DejaVu Serif}
\setsansfont{DejaVu Sans}
\usepackage{microtype}
\usepackage{xcolor}
\usepackage{hyperref}
\usepackage{enumitem}
\usepackage{parskip}
\definecolor{Accent}{HTML}{971D32}
\definecolor{Muted}{HTML}{666666}
\hypersetup{colorlinks=true,urlcolor=Accent,linkcolor=Accent}
\setlist{leftmargin=1.6em,itemsep=0.3em,topsep=0.35em}
\setcounter{tocdepth}{2}
\renewcommand{\contentsname}{Contents}
\newcommand{\sectionrule}{\par\vspace{-0.5em}\noindent\textcolor{Accent}{\rule{\linewidth}{0.6pt}}\par}
\begin{document}

\begin{center}
{\sffamily\bfseries\color{Accent} TOEFL WORD OF THE DAY\par}
\vspace{8mm}
{\fontsize{42}{48}\selectfont\bfseries ubiquitous\par}
\vspace{2mm}
{\Large\sffamily\color{Accent} /juːˈbɪkwətəs/\par}
\vspace{2mm}
{\small\sffamily Local audio phrase: ubiquitous\par}
{\small\sffamily Audio file: audios/ubiquitous.mp3\par}
\vspace{2mm}
{\large\sffamily\color{Muted} adjective\par}
\vspace{5mm}
{\sffamily present almost everywhere\par}
\vspace{6mm}
{\large Ubiquitous describes something that is so widespread that it seems to be present almost everywhere.\par}
\vspace{6mm}
\begin{minipage}{0.84\textwidth}
\itshape “Smartphones have become ubiquitous in many parts of the world.”
\end{minipage}\par
\vspace{10mm}
\textbf{Date:} August 17th 2026\quad\textbullet\quad \textbf{Level:} B1--mid-B2 TOEFL
\vfill
Created and compiled by \textbf{Amir Kooshky}\par
\vspace{2mm}
\href{https://t.me/KooshkyTOEFL}{Telegram Channel}\quad\textbullet\quad
\href{https://www.instagram.com/kooshkytoefl}{Instagram}\quad\textbullet\quad
\href{https://t.me/KooshkyTOEFL\_pv}{Personal Telegram}
\end{center}
\newpage
\tableofcontents
\newpage

\section{Meanings and usage}
\sectionrule
\subsection{Meaning 1: Present almost everywhere}
\textbf{Part of speech:} adjective

\textbf{IPA:} /juːˈbɪkwətəs/

\textbf{Pronunciation phrase:} ubiquitous

\textbf{Local audio file:} audios/ubiquitous.mp3

\textbf{Register:} formal; common in academic, technological, social, environmental, and journalistic writing

\textbf{Definition:} present, found, or appearing almost everywhere

\textbf{In simpler words:} found almost everywhere

\textbf{Typical pattern:} ubiquitous + object, technology, feature, presence, or phenomenon

\subsubsection{Natural examples}
\begin{enumerate}
  \item Smartphones have become ubiquitous in many parts of the world.
  \item Plastic is now a ubiquitous material in modern consumer products.
  \item Coffee shops are ubiquitous in the city's business districts.
  \item Digital technology has become ubiquitous in modern education.
  \item Advertising is ubiquitous on websites, social media platforms, and mobile applications.
  \item The bicycle is a ubiquitous form of transportation in some European cities.
  \item In many urban areas, security cameras have become increasingly ubiquitous.
  \item The internet has made access to information nearly ubiquitous among people with reliable digital connections.
\end{enumerate}

\subsubsection{Nuance and implication}
\textbf{Central nuance:} Ubiquitous is stronger than simply common. It suggests that something appears in so many places that it seems almost impossible to avoid.

\textbf{Usual implication:} The word often emphasizes how thoroughly something has spread through a society, environment, or area.

\textbf{Compared with rare:} Something rare is found only occasionally or in a few places. Something ubiquitous is found almost everywhere.

\subsubsection{Grammar notes}
\textbf{How to use it:} Use it before a noun, as in ubiquitous technology, or after a linking verb, as in Smartphones are ubiquitous.

\textbf{What can follow it:} It commonly describes technology, devices, materials, products, advertising, symbols, species, and features of modern life.

\textbf{Common preposition:} When saying where something is widespread, use in, throughout, or across, as in ubiquitous throughout the region.

\textbf{Position in a sentence:} Words such as increasingly, nearly, seemingly, and almost can appear before ubiquitous.

\subsubsection{Useful patterns}
\textbf{Pattern:} become ubiquitous

\textbf{Use:} Use this when something spreads until it is found almost everywhere.

\textbf{Example:} Online payment systems have become ubiquitous in many urban areas.

\textbf{Pattern:} be ubiquitous in + place or field

\textbf{Use:} Use this when something is found almost everywhere within a particular area or activity.

\textbf{Example:} Digital devices are ubiquitous in modern classrooms.

\textbf{Pattern:} ubiquitous throughout + place or system

\textbf{Use:} Use this to emphasize that something is present across nearly all parts of a larger area.

\textbf{Example:} The symbol became ubiquitous throughout the country.

\textbf{Pattern:} an ubiquitous presence

\textbf{Use:} Use this for something or someone that seems to be present almost everywhere.

\textbf{Example:} Advertising has an ubiquitous presence in modern urban life.

\textbf{Pattern:} increasingly ubiquitous

\textbf{Use:} Use this when something is becoming more widespread over time.

\textbf{Example:} Electric charging stations are becoming increasingly ubiquitous.

\subsubsection{Common wrong usage}
\paragraph{Correction 1}
\textbf{Wrong:} The device is ubiquitous only in one small laboratory.

\textbf{Correct:} The device is common only in one small laboratory.

\textbf{Why:} Ubiquitous suggests that something is present almost everywhere, so it is not suitable for something limited to one small place.

\paragraph{Correction 2}
\textbf{Wrong:} Smartphones are ubiquitously in modern society.

\textbf{Correct:} Smartphones are ubiquitous in modern society.

\textbf{Why:} Use the adjective ubiquitous after are.

\paragraph{Correction 3}
\textbf{Wrong:} The internet has became ubiquitous.

\textbf{Correct:} The internet has become ubiquitous.

\textbf{Why:} After has, use the past participle become.

\paragraph{Correction 4}
\textbf{Wrong:} The product is very ubiquitous in only a few stores.

\textbf{Correct:} The product is available in only a few stores.

\textbf{Why:} Ubiquitous conflicts with the idea of being present in only a few places.

\section{Word family}
\sectionrule
\subsection{ubiquity}
\textbf{Part of speech:} uncountable noun

\textbf{IPA:} /juːˈbɪkwəti/

\textbf{Definition:} the state of being present or found almost everywhere

\textbf{Relationship to the headword:} This noun names the condition of being ubiquitous.

\textbf{Grammar pattern:} the ubiquity of + object, technology, practice, or phenomenon

\textbf{Usage note:} It is formal and normally used without a or an.

\subsubsection{Examples}
\begin{enumerate}
  \item The ubiquity of smartphones has changed how people access information.
  \item Researchers have examined the ubiquity of plastic waste in marine environments.
\end{enumerate}

\subsection{ubiquitously}
\textbf{Part of speech:} adverb

\textbf{IPA:} /juːˈbɪkwətəsli/

\textbf{Definition:} in a way that is present or found almost everywhere

\textbf{Relationship to the headword:} This adverb describes something that exists or occurs almost everywhere.

\textbf{Grammar pattern:} ubiquitously + present, available, used, or distributed

\textbf{Usage note:} It is much less common than ubiquitous.

\subsubsection{Examples}
\begin{enumerate}
  \item Wireless devices are now used ubiquitously in many workplaces.
  \item The symbol appeared ubiquitously across the advertising campaign.
\end{enumerate}

\section{Common collocations}
\sectionrule
\subsection{become ubiquitous}
\textbf{Applies to:} Present almost everywhere

\textbf{Meaning:} spread until something is found almost everywhere

\textbf{Example:} Portable computers became ubiquitous in universities over the following decades.

\subsection{increasingly ubiquitous}
\textbf{Applies to:} Present almost everywhere

\textbf{Meaning:} becoming more widespread over time

\textbf{Example:} Artificial intelligence tools are becoming increasingly ubiquitous in professional settings.

\subsection{ubiquitous presence}
\textbf{Applies to:} Present almost everywhere

\textbf{Meaning:} the condition of appearing to be present nearly everywhere

\textbf{Example:} The ubiquitous presence of screens has changed how people consume information.

\subsection{ubiquitous technology}
\textbf{Applies to:} Present almost everywhere

\textbf{Meaning:} technology that is widely available and used almost everywhere

\textbf{Example:} Wireless communication has become a ubiquitous technology in modern society.

\subsection{ubiquitous feature}
\textbf{Applies to:} Present almost everywhere

\textbf{Meaning:} a feature found in nearly all examples of something

\textbf{Example:} Air conditioning is a ubiquitous feature of many modern office buildings.

\subsection{ubiquitous use}
\textbf{Applies to:} Present almost everywhere

\textbf{Meaning:} use that occurs almost everywhere

\textbf{Example:} The ubiquitous use of plastic packaging has created major environmental challenges.

\subsection{ubiquitous in}
\textbf{Applies to:} Present almost everywhere

\textbf{Meaning:} found almost everywhere within a particular place or field

\textbf{Pattern:} be ubiquitous in + place, society, industry, or field

\textbf{Example:} Mobile phones are ubiquitous in most large cities.

\subsection{ubiquitous throughout}
\textbf{Applies to:} Present almost everywhere

\textbf{Meaning:} present across almost all parts of a place

\textbf{Pattern:} be ubiquitous throughout + place

\textbf{Example:} The plant is ubiquitous throughout the region.

\subsection{the ubiquity of}
\textbf{Applies to:} Present almost everywhere

\textbf{Meaning:} the fact that something is present almost everywhere

\textbf{Pattern:} the ubiquity of + phenomenon

\textbf{Example:} The ubiquity of online advertising has changed the economics of media.

\textbf{Usage note:} Ubiquity is the noun form.

\section{Synonyms and antonyms}
\sectionrule
\subsection{Close synonyms}
\subsubsection{widespread}
\textbf{Part of speech:} adjective

\textbf{Applies to:} Present almost everywhere

\textbf{Shared meaning:} Both describe something found in many places.

\textbf{Difference:} Widespread means existing over a large area or among many people. Ubiquitous is stronger and suggests that something seems to be present almost everywhere.

\textbf{Example:} The new technology received widespread adoption across the industry.

\subsubsection{pervasive}
\textbf{Part of speech:} adjective

\textbf{Applies to:} Present almost everywhere

\textbf{Shared meaning:} Both describe something that is very widely present.

\textbf{Difference:} Pervasive emphasizes spreading through and affecting many parts of a system or environment. Ubiquitous emphasizes being present almost everywhere.

\textbf{Example:} Advertising has a pervasive influence on consumer behavior.

\subsubsection{omnipresent}
\textbf{Part of speech:} adjective

\textbf{Applies to:} Present almost everywhere

\textbf{Shared meaning:} Both can describe something that appears to be everywhere.

\textbf{Difference:} Omnipresent literally suggests being present everywhere and is often stronger or more literary. Ubiquitous usually means found almost everywhere in practical, real-world contexts.

\textbf{Example:} Digital screens seem almost omnipresent in major airports.

\subsubsection{prevalent}
\textbf{Part of speech:} adjective

\textbf{Applies to:} Present almost everywhere

\textbf{Shared meaning:} Both can describe something that is commonly found.

\textbf{Difference:} Prevalent means common or frequently occurring in a particular group or place. Ubiquitous emphasizes presence across nearly all places or contexts.

\textbf{Example:} The disease is particularly prevalent in tropical regions.

\subsection{Direct or useful antonyms}
\subsubsection{rare}
\textbf{Part of speech:} adjective

\textbf{Applies to:} Present almost everywhere

\textbf{Opposition:} Rare describes something found or occurring only occasionally, while ubiquitous describes something present almost everywhere.

\textbf{Example:} The species is rare and survives in only a few isolated habitats.

\subsubsection{scarce}
\textbf{Part of speech:} adjective

\textbf{Applies to:} Present almost everywhere

\textbf{Opposition:} Scarce means insufficient in amount or difficult to find, while ubiquitous means very widely present.

\textbf{Usage note:} Scarce is an opposite mainly when ubiquitous describes the availability or presence of something.

\textbf{Example:} Clean drinking water remains scarce in some drought-affected regions.

\section{Words learners may confuse}
\sectionrule
\subsection{pervasive}
\textbf{Part of speech:} adjective

\textbf{Why learners confuse them:} Both describe something that is extremely widespread, but they emphasize different aspects of that widespread presence.

\textbf{Key difference:} Ubiquitous emphasizes being present almost everywhere. Pervasive emphasizes spreading through and influencing many parts of a system or environment.

\textbf{ubiquitous:} Smartphones are ubiquitous in modern cities.

\textbf{pervasive:} Smartphones have had a pervasive influence on social communication.

\subsection{universal}
\textbf{Part of speech:} adjective

\textbf{Why learners confuse them:} Both can suggest very broad distribution, but universal focuses on inclusion of all members, not physical presence everywhere.

\textbf{Key difference:} Ubiquitous means found almost everywhere. Universal means applying to, involving, or being shared by everyone or everything in a particular group.

\textbf{ubiquitous:} Internet-connected devices are ubiquitous in many wealthy countries.

\textbf{universal:} The need for clean water is universal.

\section{Etymology}
\sectionrule
\textbf{Origin:} Ubiquitous comes from Latin ubique, meaning everywhere.

\textbf{Memory link:} The idea of everywhere is still central to the modern meaning of ubiquitous.

\section{Practice exercises}
\sectionrule
\subsection{Word family choice}
Choose the best member of the word family for each blank.

\begin{enumerate}
  \item The \_\_\_\_\_ of smartphones has transformed communication in many societies.
  \item Wireless networks have become increasingly \_\_\_\_\_ in public spaces.
  \item Digital devices are now used \_\_\_\_\_ across many industries.
\end{enumerate}

\subsection{Correct or incorrect?}
Decide whether each sentence is correct. Correct the inaccurate ones.

\begin{enumerate}
  \item Smartphones have become ubiquitous in many countries.
  \item The device is ubiquitous in one small laboratory and nowhere else.
  \item Advertising has an ubiquitous presence in many large cities.
  \item The internet is ubiquitously in modern society.
\end{enumerate}

\subsection{Collocation practice}
Complete each sentence with a natural collocation.

\begin{enumerate}
  \item Mobile payment systems have become increasingly \_\_\_\_\_ in large cities.
  \item The \_\_\_\_\_ of plastic products has created serious environmental challenges.
  \item Digital screens are ubiquitous \_\_\_\_\_ airports, shopping centers, and offices.
  \item The species is ubiquitous \_\_\_\_\_ the coastal region.
\end{enumerate}

\subsection{Complete answer key}
\subsubsection{Word family choice}
\begin{enumerate}
  \item \textbf{ubiquity.} The ubiquity of smartphones has transformed communication in many societies. \emph{Use the noun ubiquity for the state of being present almost everywhere.}
  \item \textbf{ubiquitous.} Wireless networks have become increasingly ubiquitous in public spaces. \emph{Use the adjective ubiquitous after become.}
  \item \textbf{ubiquitously.} Digital devices are now used ubiquitously across many industries. \emph{Use the adverb ubiquitously to describe how widely the devices are used.}
\end{enumerate}

\subsubsection{Correct or incorrect?}
\begin{enumerate}
  \item \textbf{Correct} \emph{Become ubiquitous is a natural pattern for something that has become present almost everywhere.}
  \item \textbf{Incorrect}. The device is found in one small laboratory and nowhere else. \emph{Ubiquitous implies very widespread presence.}
  \item \textbf{Correct} \emph{Ubiquitous presence is a natural expression for something that seems to appear almost everywhere.}
  \item \textbf{Incorrect}. The internet is ubiquitous in modern society. \emph{Use the adjective ubiquitous after is.}
\end{enumerate}

\subsubsection{Collocation practice}
\begin{enumerate}
  \item \textbf{ubiquitous.} Mobile payment systems have become increasingly ubiquitous in large cities. \emph{Increasingly ubiquitous means becoming more widely present over time.}
  \item \textbf{ubiquity.} The ubiquity of plastic products has created serious environmental challenges. \emph{The ubiquity of something means its widespread presence.}
  \item \textbf{in.} Digital screens are ubiquitous in airports, shopping centers, and offices. \emph{Use ubiquitous in when identifying places where something is widely present.}
  \item \textbf{throughout.} The species is ubiquitous throughout the coastal region. \emph{Ubiquitous throughout emphasizes presence across nearly all parts of a place.}
\end{enumerate}

\vfill
\begin{center}
\small\color{Muted} Created and compiled by \textbf{Amir Kooshky}\\
\href{https://t.me/KooshkyTOEFL}{Telegram Channel}\quad\textbullet\quad
\href{https://www.instagram.com/kooshkytoefl}{Instagram}\quad\textbullet\quad
\href{https://t.me/KooshkyTOEFL\_pv}{Personal Telegram}
\end{center}
\end{document}
